\subsection{Cell-Integrated Conservative Semi-Lagrangian Transport}
\label{sec:cell-integrated-sldg}

\begin{figure}[htbp]
\centering
\textbf{Physics flow: cell-integrated conservative semi-Lagrangian module}
\par\smallskip
\begin{tikzpicture}[node distance=0.22cm]
\node[physicsflow] (sl1) {Start from the matter stored in every spherical phase space cell};
\node[physicsflow,below=of sl1] (sl2) {Reconstruct a smooth nonnegative distribution inside each occupied source cell};
\node[physicsflow,below=of sl2] (sl3) {Compute the current spherical gravitational field and the corresponding particle trajectories};
\node[physicsflow,below=of sl3] (sl4) {Move the source-cell boundary points forward to determine where that cell's matter arrives};
\node[physicsflow,below=of sl4] (sl5) {Refine distorted maps and separate pieces that follow different reflection or central-return paths};
\node[physicsflow,below=of sl5] (sl6) {Measure the overlap of every arrival region with the destination cells and deposit the integrated matter};
\node[physicsflow,below=of sl6] (sl7) {Store unresolved central passages in the delayed-return carrier and record matter that escapes};
\node[physicsflow,below=of sl7] (sl8) {Update midpoint gravity and transport limits, check conservation and positivity, then continue};
\draw[physicsarrow] (sl1) -- (sl2);
\draw[physicsarrow] (sl2) -- (sl3);
\draw[physicsarrow] (sl3) -- (sl4);
\draw[physicsarrow] (sl4) -- (sl5);
\draw[physicsarrow] (sl5) -- (sl6);
\draw[physicsarrow] (sl6) -- (sl7);
\draw[physicsarrow] (sl7) -- (sl8);
\end{tikzpicture}
\end{figure}

\subsubsection{Purpose and distinction from nodal pullback}

The cell-integrated semi-Lagrangian discontinuous-Galerkin module is the
conservative spherical Vlasov--Poisson transport path.  In this manual,
\emph{semi-Lagrangian} means that the numerical
method follows phase space characteristics over a finite time interval.
\emph{Discontinuous Galerkin} means that the distribution is represented by a
separate polynomial inside every phase space cell; the polynomial is allowed
to jump across cell faces.  The abbreviation SLDG will be used for this
cell-integrated semi-Lagrangian discontinuous-Galerkin method.

The implemented state is a cell average
\[
 \bar f_{ijk}
 =\frac{1}{\Delta r_i\Delta w_j\Delta\eta_k}
   \int_{I_i}\!\dd r\int_{J_j}\!\dd w\int_{K_k}\!\dd\eta\,
   f(r,w,\eta),
 \qquad \eta\equiv\frac{L^2}{2}.
\]
Here \(I_i=[r_{i-1/2},r_{i+1/2}]\), \(J_j=[w_{j-1/2},w_{j+1/2}]\),
and \(K_k=[\eta_{k-1/2},\eta_{k+1/2}]\) are one radial, radial-velocity,
and angular momentum cell.  The symbols
\(\Delta r_i\), \(\Delta w_j\), and \(\Delta\eta_k\) denote their widths.
Since \(L\,\dd L=\dd\eta\), the physical phase space measure is
\[
 \dd\Gamma=8\pi^2\,\dd r\,\dd w\,\dd\eta,
 \qquad
 M_{ijk}=8\pi^2\bar f_{ijk}
          \Delta r_i\Delta w_j\Delta\eta_k .
\]
The factor \(8\pi^2\) results from integrating the two configuration-space
angles and the unused velocity azimuth under spherical symmetry.

The nodal semi-Lagrangian method traces each target grid point back to its
departure point and interpolates \(f\) there.  The SLDG method transports
integrated mass: it reconstructs every \emph{source cell}, maps that cell
forward, and integrates its overlap with the destination cells.  This
source-to-destination accounting gives the conservative update
\[
 M^{n+1}_{abc}
 =\sum_{ijk}
   \int_{\mathcal A_{ijk}\cap D_{abc}}
   f^n_{h,ijk}\circ\mathcal X^{-1}\,
   8\pi^2\,\dd r\,\dd w\,\dd\eta ,
\]
where \(\mathcal X\) is the characteristic map from time \(t^n\) to
\(t^{n+1}\), \(\mathcal A_{ijk}=\mathcal X(I_i\times J_j\times K_k)\)
is an arrival region, \(D_{abc}\) is a destination cell, and
\(f^n_{h,ijk}\) is the reconstructed source polynomial.  There is no
post-step global mass rescaling.  Source mass outside the visible grid is
recorded either in the delayed-inner carrier, which stores unresolved central
passages, or in the escape budget.  Together these destinations account for
the complete transported mass.

The relevant sources are:
\begin{itemize}
\item \path{cpu/include/ev/conservative_sldg_transport.hpp}: polynomial
reconstruction, mapped-polygon geometry, recursive refinement, exact overlap
integration, delayed-return coupling, and midpoint-field driver;
\item \path{cpu/include/ev/hamiltonian_sldg_transport.hpp}: Hamiltonian-aligned
reconstruction, cached modal state, and static-potential timestep control;
\item \path{cpu/include/ev/hamiltonian_orbit_reference.hpp}: optional positive
orbit-average reference \(f_0(H,\eta)\);
\item \path{cpu/include/ev/delayed_inner_orbit_carrier.hpp}: persistent
finite-volume mass inside \(0<r<r_{\min}\);
\item \path{cpu/include/ev/radial_orbit_regularization.hpp}: selective
Levi--Civita characteristic retry near stiff pericentres.
\end{itemize}

\subsubsection{Frozen spherical Hamiltonian map}

For static spherical Vlasov--Poisson evolution at fixed
\(\eta=L^2/2\), the one-particle Hamiltonian is
\[
 H(r,w;\eta)
 =\frac{w^2}{2}+\frac{\eta}{r^2}+\Phi(r).
\]
Here \(w=\dot r\) is radial velocity, \(L\) is the magnitude of specific
angular momentum, \(\Phi\) is the gravitational potential, and
\(\eta/r^2=L^2/(2r^2)\) is tangential kinetic energy.  The characteristic
equations are
\[
 \dot r = \frac{\partial H}{\partial w}=w,
 \qquad
 \dot w =-\frac{\partial H}{\partial r}
 =\frac{2\eta}{r^3}-\Phi_r,
 \qquad
 \dot\eta=0,
\]
where \(\Phi_r=\partial\Phi/\partial r\).  Hamiltonian flow preserves
phase space area:
\[
 \det\!\left(\frac{\partial(r',w')}{\partial(r,w)}\right)=1,
 \qquad
 \dd r'\,\dd w'=\dd r\,\dd w .
\]
This area preservation is both a physical property and a numerical geometry
check.  For a small source patch with area \(A_0\), the mapped corner polygon
should have area \(A\simeq A_0\).  VENUS records the relative area defect
\[
 \epsilon_A=\frac{|A-A_0|}{A_0}.
\]

The default characteristic update inside the SLDG map is
kick--drift--kick (KDK).  If
\[
 a_r(r;\eta) = \frac{2\eta}{r^3}-\Phi_r(r)
\]
denotes radial acceleration in the frozen field, one KDK step is
\[
 w^{n+1/2}=w^n+\frac{\Delta t}{2}a_r(r^n;\eta),
\quad
 r^{n+1}=r^n+\Delta t\,w^{n+1/2},
\quad
 w^{n+1}=w^{n+1/2}
 +\frac{\Delta t}{2}a_r(r^{n+1};\eta).
\]
The exact implementation also resolves the selected radial boundary and may
retry a stiff pericentre segment with the Levi--Civita transformation described
in the orbit-regularization module reference below.  The frozen-field map is used
only within one transport evaluation; self-gravity is time-centred as
described below.

\subsubsection{P2 source polynomial and its cell average}

Inside one \((r,w)\) source cell at a fixed \(\eta\) quadrature node, define
normalized coordinates
\[
 x=\frac{r-r_i}{\Delta r_i},
 \qquad
 y=\frac{w-w_j}{\Delta w_j},
 \qquad -\frac12\le x,y\le\frac12 .
\]
The notation \(P_2\) means the six-dimensional space of polynomials in
\((x,y)\) whose total degree is at most two: one constant, two linear, and
three quadratic modes.  The implemented \(P_2\) polynomial is
\begin{equation*}
 \begin{split}
 f_h(x,y)={}&\bar f
 +a_r x+a_w y
 +a_{rr}\left(x^2-\frac1{12}\right)\\
 &+a_{rw}xy
 +a_{ww}\left(y^2-\frac1{12}\right).
 \end{split}
 \tag{SLDG-P2}\label{eq:sldg-p2}
\end{equation*}
All symbols in this expression are local to one source cell:
\(\bar f\) is the stored cell average; \(a_r,a_w\) are linear modes;
\(a_{rr},a_{rw},a_{ww}\) are quadratic modes.  Since
\[
 \int_{-1/2}^{1/2}x\,\dd x=0,
 \qquad
 \int_{-1/2}^{1/2}\left(x^2-\frac1{12}\right)\dd x=0,
\]
and similarly for \(y\), every nonconstant mode has zero mean.  Therefore the
polynomial integrates to \(\bar f\Delta r_i\Delta w_j\) independently of
the slope and curvature values.

The plain conservative SLDG mode currently constructs limited linear slopes.
For a cell-centred sequence \(q_i\), the minmod limiter is
\[
 \operatorname{minmod}(a,b,c)=
 \begin{cases}
  \min(a,b,c),&a,b,c>0,\\
  \max(a,b,c),&a,b,c<0,\\
  0,&\text{otherwise},
 \end{cases}
\]
and the radial slope is obtained from the backward, centred, and forward
differences with the usual factor-of-two bounds.  The same construction is
used in \(w\). At a wall of fixed physical radius \(R\), the reflecting
radial ghost state is evaluated at \(w'=-w-2HR\), where \(H\) is the
Hubble rate. The static limit \(H=0\) gives \(w'=-w\).

The source representation in
\path{cpu/include/ev/conservative_sldg_transport.hpp} is:
\codesource{Conservative SLDG source polynomial}
{\path{cpu/include/ev/conservative_sldg_transport.hpp}}
{Stores the cell average and five local quadratic modes and applies the
 source-cell positivity limiter before forward transport}
\begin{lstlisting}[language=C++]
struct SLDGSourcePolynomial {
  double cell_average = 0.0;
  std::array<double, 5> mode{};
  double center_r = 0.0, center_w = 0.0;
  double scale_r = 1.0, scale_w = 1.0;
  double enforce_positivity();
};
\end{lstlisting}
The five entries of \code{mode} store the slopes and curvatures
\((a_r,a_w,a_{rr},a_{rw},a_{ww})\) in
Equation~\eqref{eq:sldg-p2}.  Each associated basis function has zero cell
average, so changing these entries preserves the stored mass.

\subsubsection{Positivity limiter}

A phase space density must satisfy \(f\ge0\).  Let
\[
 f_h=\bar f+\delta f,
\qquad \langle\delta f\rangle_{\rm cell}=0,
\]
and let \(f_{\min}\) be the minimum of the reconstructed polynomial over the
cell.  If \(f_{\min}<0\), VENUS scales every nonconstant mode by one factor:
\[
 f_h^{\rm lim}=\bar f+\theta\delta f,
\qquad
 \theta=\min\!\left(1,
 \frac{\bar f}{\bar f-f_{\min}}\right).
\]
Thus \(0\le\theta\le1\), the cell average is unchanged, and the limited
minimum is nonnegative.  Because a quadratic can attain its minimum inside an
edge or the cell interior, the implementation checks the relevant stationary
points and uses a robust bisection correction if roundoff prevents the direct
factor from satisfying the positivity criterion.  The limiter is local: it does not
move mass to a neighbouring cell.

\subsubsection{Subcells, trace lattice, and map-quality checks}

Each \((r,w)\) source cell is divided into
\(\code{vp_sldg_subcells}\) pieces in each direction.  The default is two,
so there are four initial patches.  Each \(\eta\) cell is integrated with
\(\code{vp_sldg_eta_quadrature_points}\) Gaussian nodes, also two by default.
At one angular node \(\eta_q\), the orbit uses
\[
 L_q=\sqrt{2\eta_q}.
\]

For one source patch VENUS traces a \(3\times3\) lattice: four corners, four
edge midpoints, and the centre.  The corner image defines an arrival
quadrilateral.  The edge-midpoint and centre images measure how far the true
characteristic map departs from a bilinear map.  If
\(\bm X=(r',w')\) denotes an exactly traced lattice point and
\(\bm X_{\rm bil}\) is the value predicted from the four mapped corners, the
normalized map defect is schematically
\[
 \epsilon_{\rm map}
 =\max_{\rm midpoints,centre}
 \left[
 \left(\frac{r'-r'_{\rm bil}}{\Delta r_{\rm map}}\right)^2+
 \left(\frac{w'-w'_{\rm bil}}{\Delta w_{\rm map}}\right)^2
 \right]^{1/2}.
\]
The mapped widths \(\Delta r_{\rm map}\) and \(\Delta w_{\rm map}\) prevent a
large-coordinate but tiny-patch map from being judged by an absolute error.

A patch is accepted only when all of the following hold:
\begin{itemize}
\item mapped coordinates are finite;
\item the corner polygon is counter-clockwise and convex;
\item all eight centre/edge/corner triangles have positive orientation;
\item the mapped area is larger than the geometry floor;
\item \(\epsilon_A\le\code{vp_sldg_max_area_defect}\);
\item \(\epsilon_{\rm map}\le\code{vp_sldg_max_map_defect}\).
\end{itemize}
The default area and map bounds are \(10^{-4}\) and \(10^{-2}\),
respectively.  The geometry floor is controlled by
\code{vp_sldg_geometry_tolerance}.

Adjacent child patches reuse parent and neighbour trace points.  This trace
cache is important because orbit integration is much more expensive than
polygon clipping.  Different \(\eta\) planes are independent and form the
outer parallel work units; a plane is owned by only one worker at a time, so
destination accumulation is race free.

\subsubsection{Recursive refinement and reflecting-wall separatrix}

An invalid mapped patch is not immediately deposited by points.  VENUS splits
it in \(r\) or \(w\), choosing the direction with the larger normalized map
curvature, and repeats the trace/quality check.  The maximum depth and work are
\code{vp_sldg_max_refinement} and
\code{vp_sldg_max_work_per_subcell}.  Their defaults are 30 and 131072.

A reflecting inner wall introduces a discontinuous branch in the map: one
part of a source patch reaches \(r_{\min}\) and changes \(w\to-w\), while a
nearby part follows the free orbit. After the gravitational half kick, let
\(w_h=w-\Delta t\,\Phi_r/2\) and \(v_t=L/r\) be the radial and
tangential velocities. The closest approach of the free angular drift occurs at
\[
 t_* = \operatorname{clip}\!\left(
 -\frac{rw_h}{w_h^2+v_t^2},
 \min(0,\Delta t),\max(0,\Delta t)\right).
\]
Here clipping restricts the time to the traced interval; a stationary orbit
uses \(t_*=0\). The boundary between wall-crossing and free trajectories is
\[
 S(r,w)=\left[(r+t_*w_h)^2+(t_*v_t)^2\right]^{1/2}-r_{\min}=0.
\]
Using the closest approach detects a wall encounter even when the
unreflected endpoint has already returned to larger radius.
When a patch straddles \(S=0\), the splitter chooses the coordinate along
which \(S\) varies most.  The wall branch receives additional permitted
refinement depth because its map is physically continuous only after the
reflection topology is resolved.
\codesource{Inner-wall refinement criterion}
{\path{cpu/include/ev/conservative_sldg_transport.hpp}}
{Evaluate the closest approach during the free angular drift before choosing a refinement direction}
\begin{lstlisting}[language=C++]
const double closest_time = speed2 > 0.0
    ? std::max(std::min(0.0, dt),
        std::min(std::max(0.0, dt), -source_r * w_half / speed2))
    : 0.0;
return std::hypot(source_r + closest_time * w_half,
    closest_time * tangential) - p.r_min;
\end{lstlisting}

Only after the depth or work budget is exhausted may VENUS use a bounded
\(2\times2\) Gauss-point deposit.  If \(M_{\rm fb}\) is the physical mass
carried by this fallback and \(M_{\rm src}\) is total source mass, the run
requires
\[
 \frac{M_{\rm fb}}{M_{\rm src}}
 \le\code{vp_sldg_fallback_mass_tolerance}.
\]
The default tolerance is \(10^{-9}\).  The diagnostic records the largest
fallback cell, its \((r,w,L)\), refinement depth, and work, so fallback cannot
silently become the production algorithm.

\subsubsection{Exact polygon overlap and moment integration}

An accepted quadrilateral is divided into eight affine triangles using its
mapped corners, edge midpoints, and centre.  Each triangle is clipped against
candidate destination rectangles by successive half-plane clipping.  Polygon
area and normalized moments are then evaluated analytically.  On each
triangle, the affine numerical map \(\mathcal X_h\) has constant inverse
Jacobian
\[
 J_{\rm inv}=\left|\det D\mathcal X_h^{-1}\right|
 =\frac{A_{\rm source}}{A_{\rm mapped}}.
\]
Here \(A_{\rm source}\) and \(A_{\rm mapped}\) are the positive areas of
that triangle before and after mapping.  The exact Hamiltonian map has
unit Jacobian; the approximate triangle map retains its measured area
change.  Conservation therefore deposits
\(J_{\rm inv} f_h\circ\mathcal X_h^{-1}\) on the arrival triangle.
For a clipped polygon \(P\), VENUS needs
\[
 A_P=\int_P\dd r\,\dd w,\quad
 R_P=\int_P x\,\dd r\,\dd w,\quad
 W_P=\int_P y\,\dd r\,\dd w,
\]
\[
 R^{(2)}_P=\int_P\left(x^2-\frac1{12}\right)\dd r\,\dd w,
\quad
 RW_P=\int_Pxy\,\dd r\,\dd w,
\quad
 W^{(2)}_P=\int_P\left(y^2-\frac1{12}\right)\dd r\,\dd w .
\]
In these moments, \(x,y\) are the normalized source-cell coordinates
evaluated through the inverse affine map.  The deposited reduced mass is
\[
 m_P=J_{\rm inv}\left[\bar f A_P+a_rR_P+a_wW_P
 +a_{rr}R^{(2)}_P+a_{rw}RW_P+a_{ww}W^{(2)}_P\right].
\]
Multiplication by \(8\pi^2\Delta\eta_k\) and the angular quadrature weight
gives physical mass.

The triangle-area ratio is applied locally before accumulation:
\codesource{SLDG conservative affine-triangle deposition}
{\path{cpu/include/ev/conservative_sldg_transport.hpp}}
{Converts arrival-polygon integrals to source mass with the inverse map Jacobian}
\begin{lstlisting}[language=C++]
const long double inverse_jacobian =
    source_triangle_area / mapped_triangle_area;
\end{lstlisting}
The destination contribution then multiplies the polynomial integral by
\code{inverse_jacobian}.  The sum over all destination intersections
recovers that source triangle's mass even when its allowed map-area defect
is nonzero.

Deep refinement can make a polygon tiny relative to its absolute coordinates.
The clipping context therefore translates and rescales each mapped polygon to
an order-unity box and performs moment arithmetic in extended precision.  This
reduces cancellation without altering the physical overlap.

\subsubsection{Delayed inner return as a finite-volume carrier}
\label{sec:sldg-delayed-carrier}

A positive logarithmic grid does not resolve \(0<r<r_{\min}\).  Treating
\(r_{\min}\) as a reflecting wall changes the orbital period of a low-angular-
momentum orbit.  The option
\code{inner_boundary = delayed_return} instead stores the mass that crosses
the inner face in a persistent carrier.

One carrier packet stores reduced mass
\[
 m_{\rm red}=\int f\,\dd r\,\dd w
\]
for one \(\eta\) cell, its actual quadrature value \(\eta_q\), boundary speed,
and remaining hidden time.  Its physical mass is
\[
 m_{\rm phys}=8\pi^2\Delta\eta_k\,\omega_q\,m_{\rm red},
\]
where \(\omega_q\) is the Gaussian weight normalized to the \(\eta\)-cell
integral.

When a mapped patch contains both visible and hidden trajectories, recursive
refinement first tries to resolve their separatrix.  If the configured carrier
refinement depth is reached, Gaussian quadrature partitions the reconstructed
mass between visible deposition and hidden packets without renormalizing
either part.  A hidden packet advances its clock; when the clock expires it is
released at \(r_{\min}\) with the outgoing radial velocity.

The transport ledger is
\[
 M_{\rm grid}^{n}+M_{\rm hidden}^{n}
 =M_{\rm grid}^{n+1}+M_{\rm hidden}^{n+1}+M_{\rm escaped}.
\]
The residual of this equation is checked against
\code{vp_inner_carrier_mass_tolerance}.  Optional return-time-bin compression
reduces a large packet list while preserving positive mass and the monitored
mass, momentum, energy, and return-time moments.

\subsubsection{Midpoint self-gravity from the continuity equation}

The final transport should feel a field representative of the middle of the
step.  VENUS predicts shell mass at \(t^{n+1/2}\) from the spherical
continuity equation
\[
 \frac{\partial\rho}{\partial t}
 +\frac{1}{r^2}\frac{\partial}{\partial r}(r^2j_r)=0,
\]
where \(j_r\) is radial mass current.  If
\[
 M_i=4\pi r_i^2\rho_i\Delta r_i,
 \qquad
 F_{i+1/2}=4\pi r_{i+1/2}^2j_{r,i+1/2},
\]
then the finite-volume predictor is
\[
 \frac{\dd M_i}{\dd t}
 =-\left(F_{i+1/2}-F_{i-1/2}\right).
\]
The face current is interpolated from cell values.  A donor limiter scales
outgoing fluxes if a half step would remove more mass than the donor shell
contains.  Inner-carrier exchange and outer escape are included in the
predictor ledger.  Poisson's equation is solved from the predicted midpoint
mass, and the full SLDG characteristic map is then evaluated through that
midpoint field.

This procedure time-centres gravity without depositing a provisional full
SLDG solution.  Diagnostics report predictor mass before and after, predicted
escape, hidden-mass change, and the number of limited faces.

\subsubsection{Hamiltonian-aligned reconstruction}

For an exact stationary spherical equilibrium,
\[
 f=f(H,\eta),
\]
so the distribution is constant along a frozen orbit.  A polynomial aligned
only with \(r\) and \(w\) can interpret this curved orbit as a large local
gradient and create unnecessary diffusion.  The Hamiltonian SLDG option uses
two local directions:
\begin{itemize}
\item an across-orbit direction measured by the change in \(H\);
\item an along-orbit coordinate \(s_{\rm orb}\), tangent to the Hamiltonian
      flow.
\end{itemize}
The coordinate \(s_{\rm orb}\) is a local, time-like parameter along the
orbit, not the Schwarzschild time \(\tau\) and not necessarily a global angle.
Its purpose is to distinguish a physical variation along an orbit from a
variation between neighbouring energy surfaces.

The cell moments
\(\langle H\rangle\) and \(\langle H^2\rangle\) are evaluated before the fit.
The velocity moments used analytically are
\[
 \langle w^2\rangle=w_j^2+\frac{\Delta w_j^2}{12},
\qquad
 \langle w^4\rangle=w_j^4+
 \frac{w_j^2\Delta w_j^2}{2}+\frac{\Delta w_j^4}{80}.
\]
The \(r\)-dependent potential and centrifugal terms use four-point Gaussian
quadrature.  Neighbouring cell averages then determine a quadratic local fit
in Hamiltonian and phase features.  If the local design matrix loses rank, or
the phase direction is degenerate near a turning point, the implementation
falls back to a lower-information reconstruction and increments an explicit
diagnostic counter.  ``Loses rank'' means that the sampled neighbouring
features no longer provide independent constraints, so more than one set of
polynomial coefficients fits the same data and the full quadratic inversion
is not trustworthy.  The fitted result is converted to the same five P2 modes
in Equation~\eqref{eq:sldg-p2}, limited for positivity, and transported by the
common geometric SLDG engine.

\subsubsection{Positive orbit-average reference}

With
\[
 \dd r\,\dd w=\dd H\,\dd s_{\rm orb}
\]
at fixed \(\eta\), the orbit average
\[
 f_0(H,\eta)
 =\frac{\int_{\rm orbit}f(H,s_{\rm orb},\eta)\,\dd s_{\rm orb}}
        {\int_{\rm orbit}\dd s_{\rm orb}}
\]
is stationary in a frozen potential.  If
\code{vp_hamiltonian_reference = orbit_average}, VENUS builds a nonnegative
binned approximation to \(f_0\), splits
\[
 f=f_0+\delta f,
\]
transports only the signed residual \(\delta f\), and adds \(f_0\) back at the
endpoint.  The reference uses closed-orbit energy ranges and is normalized to the
represented mass. Its amplitude is reduced when needed to keep the total
reconstructed distribution nonnegative.  Empty energy bins are filled only from
physically admissible neighbouring closed-orbit bins.

This is a numerical reference state, not extra matter.  Its value must not be
added to Poisson separately.  Diagnostics report the reference mass fraction,
the residual \(L^1\) fraction, active \(\eta\) planes, the minimum positivity
scale, and how many planes required scaling.  Collisions change the modal
distribution and therefore invalidate the cached Hamiltonian reconstruction.
Versioned grid restarts preserve the five modal coefficients so a clean
restart does not rebuild them from cell averages.

\subsubsection{Adaptive transport cap}

If a positive initial SLDG transport cap is configured, the next cap responds
to the measured fallback fraction.  Let
\[
 \epsilon_{\rm fb}=\frac{M_{\rm fb}}{M_{\rm src}},
 \qquad
 \epsilon_{\rm target}
 =\code{vp_sldg_target_fallback_fraction}\,
  \code{vp_sldg_fallback_mass_tolerance}.
\]
The raw scale factor is
\[
 g=\sqrt{\frac{\epsilon_{\rm target}}{\epsilon_{\rm fb}}},
\]
with appropriate growth behavior when \(\epsilon_{\rm fb}=0\).  It is clipped
between \code{vp_sldg_timestep_shrink} and
\code{vp_sldg_timestep_growth}.  This cap is combined with the physical CFL,
dynamical, potential-variation, output, and final-time limits; it does not
replace them.

For Hamiltonian SLDG, the static-field timescale also estimates
\[
 t_\Phi=\min_i\left(
 \left|\frac{\Phi_i}{\dot\Phi_i}\right|,
 \left|\frac{g_i}{\dot g_i}\right|
 \right),
\]
over reliable cells, where \(g_i=-\Phi_{r,i}\).  The macrostep is bounded by
\(\code{vp_potential_variation_eta}\,t_\Phi\).

\subsubsection{How to run the present SLDG implementation}

Save the following small conservative setup as \path{params.ini}.  The
\code{kinetic_grid} and \code{vp_interpolation} lines are the essential SLDG
selectors.
\codesource{\code{venus_cpu_vlasov_poisson_eulerian}}
{\path{cpu/src/apps/vlasov_poisson/eulerian.cpp};
 \path{cpu/include/ev/config.hpp}; \path{cpu/include/ev/solver.hpp}}
{Parses the conservative cell-integrated transport controls under the locked
 collisionless Vlasov--Poisson task profile}
\begin{lstlisting}[basicstyle=\ttfamily\scriptsize]
# params.ini: conservative cell-integrated SLDG run
physics_model = vlasov_poisson
vp_universe = static
kinetic_grid = cell_average_eta
vp_interpolation = conservative_sldg
initial_data = analytic
shell_initial_df = true
gaussian_density_shell = true
Nr = 24
Nw = 24
NL = 12
r_min = 0.1
r_max = 2
w_max = 2
L_max = 0.2
integrator = semi_lagrangian_kick_drift_kick
vp_timestep_mode = frank_dynamical
vp_dynamical_eta = 0.02
vp_sldg_subcells = 2
vp_sldg_fallback_mass_tolerance = 1e-8
inner_boundary = delayed_return
t_end = 0.04
dt_max = 0.01
output_dir = output/vp_sldg_example
\end{lstlisting}
Build and run the task-specific Vlasov--Poisson frontend:
\codesource{\code{venus_cpu_vlasov_poisson_eulerian}}
{\path{cpu/src/apps/vlasov_poisson/eulerian.cpp}; \path{cpu/Makefile}}
{Builds and launches the locked collisionless frontend with the supplied
 semi-Lagrangian configuration}
\begin{lstlisting}
cd VENUS/cpu
make venus_cpu_vlasov_poisson_eulerian

./venus_cpu_vlasov_poisson_eulerian --config params.ini
\end{lstlisting}
For a new configuration, the essential conservative SLDG block is:
\codesource{Conservative SLDG configuration}
{\path{cpu/include/ev/config.hpp};
 \path{cpu/include/ev/conservative_sldg_transport.hpp}}
{Selects the cell-average geometry, conservative forward map, refinement
 checks, and fail-closed fallback-mass bound}
\begin{lstlisting}
physics_model = vlasov_poisson
kinetic_grid = cell_average_eta
vp_interpolation = conservative_sldg
integrator = semi_lagrangian_kick_drift_kick
vp_sldg_subcells = 2
vp_sldg_eta_quadrature_points = 2
vp_sldg_max_area_defect = 1e-4
vp_sldg_max_map_defect = 1e-2
vp_sldg_fallback_mass_tolerance = 1e-9
\end{lstlisting}
To select the Hamiltonian-aligned reconstruction:
\codesource{Hamiltonian-aligned SLDG configuration}
{\path{cpu/include/ev/hamiltonian_sldg_transport.hpp};
 \path{cpu/include/ev/hamiltonian_orbit_reference.hpp}}
{Reconstructs along frozen Hamiltonian geometry and optionally supplies a
 positive orbit-average reference distribution}
\begin{lstlisting}
vp_universe = static
vp_interpolation = hamiltonian_sldg
vp_hamiltonian_reference = orbit_average
vp_hamiltonian_reference_energy_bins = 512
vp_hamiltonian_reference_quadrature_points = 4
\end{lstlisting}
To retain the physical delay through the unresolved central region:
\codesource{SLDG delayed-inner-return carrier}
{\path{cpu/include/ev/conservative_sldg_transport.hpp};
 \path{cpu/include/ev/delayed_inner_orbit_carrier.hpp}}
{Carries unresolved central trajectories for their computed return time while
 enforcing quadrature, refinement, segment, and mass tolerances}
\begin{lstlisting}
inner_boundary = delayed_return
vp_inner_carrier_quadrature_points = 4
vp_inner_carrier_max_refinement = 6
vp_inner_carrier_max_orbit_segments = 16
vp_inner_carrier_mass_tolerance = 1e-12
\end{lstlisting}

The parameter consistency checks require cell-average \(\eta\) geometry for SLDG.
Hamiltonian SLDG and its orbit reference require a static universe.  Prescribed
outer inflow is not part of the present cell-integrated Hamiltonian path.
\code{delayed_return} is accepted only with
\code{conservative_sldg} or \code{hamiltonian_sldg}.

\subsubsection{Implemented algorithms and applicability limits}

The implementation provides:
\begin{itemize}
\item source-cell conservative forward transport with no global mass fix;
\item positivity-limited local polynomial reconstruction;
\item cached trace lattices, geometric validity checks, recursive refinement,
and exact polygon/rectangle overlap moments;
\item explicit, tolerance-gated fallback ledger and adaptive transport cap;
\item midpoint self-gravity from a conservative continuity predictor;
\item reflecting, outflow, and delayed-inner-return branches;
\item persistent hidden-mass carrier with optional conservative compression;
\item Hamiltonian-aligned P2 reconstruction and an optional positive
orbit-average reference;
\item OpenMP/std::thread work over independent \(\eta\) planes and MPI
\(\eta\)-slab reductions;
\item restart support for the grid, carrier, adaptive state, and Hamiltonian
modal coefficients;
\item diagnostics for every mass destination and every map-quality fallback.
\end{itemize}

The following limitations govern scientific use:
\begin{itemize}
\item the Hamiltonian fit and orbit-reference parameters must be refined for
equilibrium and strongly time-dependent problems;
\item fallback, rank-deficiency, phase-degeneracy, and positivity diagnostics
must remain acceptance gates rather than merely logged numbers;
\item delayed-return separatrices and carrier compression require special
care for very low \(L\), long hidden periods, and field changes during the
delay;
\item expanding-universe cell-integrated Hamiltonian transport and prescribed
outer inflow are not yet the supported Hamiltonian-SLDG production path;
\item collisions invalidate cached transport modes, so a coupled calculation
must close its transport and collision mass ledgers and preserve restart
equivalence.
\end{itemize}

The conservative ledgers, midpoint coupling, and Hamiltonian reconstruction
make every fallback and unresolved destination explicit.  Scientific use
requires problem-specific refinement until fallback mass, carrier mass, and
changes in the reported observables are below the chosen error budget.
